% Anya
% biological-perspectives-coevolution.tex
%%%%%%%%%%%%%%%%%%%%%%%%%%%%%%%%%%%%%%%%%%%%%%%%%%%%%%%%%%%%%%%%%%%%%%%%%%%%%%%%%%
% revision status as of 3/10: x? 
%           There is certainly lots more that could be added, but I don't think 
%           much more can be without a much more involved literature review of the 
%           EC co-evolutionary literature, which I (Anya) can't do at this point
% - [x] accept or reject all track changes
% - [x] incorporate review comments
%    - don’t worry about clearing all comments
% - [x] incorporate reference suggestions
% - [x] cross-cutting revision objectives (see email)
%     - more references to EC and call backs added
% - [x] other planned changes or refinements (if any)
% for revision status, mark - for partial or x for complete
% any free-form comments on planned/completed revisions are helpful too!
% The comments that remain are ones that are bigger lifts and so I don't disagree with them, but don't feel equipped to write that additional text at this point. There is so much more that could be said and analyzed, but it would be its own project at that point I think!
% --------------------------------------------------------------------------------
% draft status: Complete Draft as of 11/10/25    <---------- fill me in!
% pick one of: outline, partial draft, rough draft, complete draft
% where rough draft indicates that major points are hit,
% but still working on changes or improvements in refining the text
%%%%%%%%%%%%%%%%%%%%%%%%%%%%%%%%%%%%%%%%%%%%%%%%%%%%%%%%%%%%%%%%%%%%%%%%%%%%%%%%%%

\section{Biological Perspectives on Coevolution}
\label{sec:coevolution}

% \textbf{Suggested Citations}: \begin{itemize}
% \item find full reference info in \texttt{filecontents} above, or in rendered PDF
% \item empty this list by moving citations into ``Included Citations'' or ``Excluded Citations'' below

% \end{itemize}

% \textbf{Included Citations}: \begin{itemize}
% \item move citations incorporated into section text here!
% \item \citet{Hillis_1990}
% \item \citet{Wang2019}
% \item \citet{Bucci2003,Bucci2005,Bucci2003a}
% \item \citet{Watson2003}
% \item \citet{Watson2001}
% \item \citet{Stanley2004}
% \item \citet{Popovici2012}
% \item \citet{Ficici2000}
% \item \citet{Cliff1995}
% \item \citet{Zaman2014}
% \item \citet{Vostinar2019}
% \end{itemize}

% \textbf{Excluded Citations}: \begin{itemize}
% \item move citations not incorporated into section text here!
% \end{itemize}

% 8. Biological Perspectives on Coevolution (Anya)
% might compare/contrast with “cooperative/competitive" coevolution” framing in EC
% wherever possible, compare/contrast with existing EC perspectives
%         - draft outline includes some possible angles/refs in some sections
% -  as primary focus, opportunities to understand existing practice
% - section names are broad — focusing on a handful of selected topics within a section is fine!

% TODO @mmore500 high-profile \cite{Hillis1990}
% pulled from citations to hillis, 1990
% 1. Potter, M. A., & De Jong, K. A. (1994). A Cooperative Coevolutionary Approach to Function Optimization. Parallel Problem Solving from Nature (PPSN III). https://doi.org/10.1007/3-540-58484-6_269

% 2. Rosin, C. D., & Belew, R. K. (1997). New Methods for Competitive Coevolution. Evolutionary Computation, 5(1), 1–29. https://doi.org/10.1162/evco.1997.5.1.1

% 4. Arcuri, A., & Yao, X. (2008). A Novel Co-evolutionary Approach to Automatic Software Bug Fixing. 2008 IEEE Congress on Evolutionary Computation (IEEE World Congress on Computational Intelligence). https://doi.org/10.1109/CEC.2008.4630793

% 5. Paredis, J. (1995). Coevolutionary Computation. Artificial Life, 2(4), 355–375. https://doi.org/10.1162/artl.1995.2.4.355

% 6. Wolpert, D. H., & Macready, W. G. (2005). Coevolutionary Free Lunches. IEEE Transactions on Evolutionary Computation, 9(6), 721–735. https://doi.org/10.1109/TEVC.2005.856205

% 8. Paredis, J. (1994). Co-evolutionary Constraint Satisfaction. Parallel Problem Solving from Nature (PPSN III). https://doi.org/10.1007/3-540-58484-6_249

% 12. Pagie, L., & Hogeweg, P. (1997). Evolutionary Consequences of Coevolving Targets. Evolutionary Computation, 5(4), 401–418. https://doi.org/10.1162/evco.1997.5.4.401

% 13. Schmidt, M. D., & Lipson, H. (2008). Coevolution of Fitness Predictors. IEEE Transactions on Evolutionary Computation, 12(6), 736–749. https://doi.org/10.1109/TEVC.2008.919006

% 14. Ficici, S. G., & Pollack, J. B. (2001). Pareto Optimality in Coevolutionary Learning. European Conference on Artificial Life (ECAL 2001). https://doi.org/10.1007/3-540-44811-X_34

% 18. Iorio, A. W., & Li, X. (2004). A Cooperative Coevolutionary Multiobjective Algorithm Using Non-dominated Sorting. Genetic and Evolutionary Computation Conference (GECCO 2004). https://doi.org/10.1007/978-3-540-24854-5_56

% 19. Panait, L., Luke, S., & Wiegand, R. P. (2006). Biasing Coevolutionary Search for Optimal Multiagent Behaviors. IEEE Transactions on Evolutionary Computation, 10(6), 629–645. https://doi.org/10.1109/TEVC.2006.880330

% also high-profile
% Paired Open-Ended Trailblazer (POET): Endlessly Generating Increasingly Complex and Diverse Learning Environments and Their Solutions Rui Wang, Joel Lehman, Jeff Clune, Kenneth O. Stanley https://doi.org/10.48550/arXiv.1901.01753


% - sections: at most ~2-3 pages (~<1k-1.5k words; page est. excludes refs)
All biological evolution is, in fact, co-evolution between many interacting species. The interaction between the species can be negative, such as parasitism, predation, or general competition, or positive, such as mutualism or cooperation. Such interactions often vary in their strength and even the very nature of the effect over evolutionary time. 
Evolutionary computation has drawn inspiration from co-evolutionary dynamics such as predation \citep{Willkens2022,Popovici2012, Ficici2000}, parasitism \citep{Hilllis_1990,  Cliff1995}, competition \citep{Stanley2004}, and mutualism \citep{Wang2019, Bucci2005, Bucci2003a, Bucci2003, Watson2003, Watson2001}, however the specific form of co-evolutionary relationship plays an often underappreciated role in these algorithms.

For example, some coevolutionary algorithms use a form of parasitism where a population of tests is evaluated against a population of solutions. %TODO: find example citations
While these are sometimes presented as a form of predation, failing a test doesn't ``kill'' a solution. Instead, failing a test leads to a slightly lower fitness score for the solution, similar to a parasitism dynamic. 

%TODO: look for places to point out that bio suggests type of EC to use or when parasitism/mutualism would change so making sure to consider if parasitism/mutualism
The classic example of co-evolutionary algorithm is \citet{Potter2000}.
In this work, they introduced the possibility of decomposing a problem into subcomponents and using separate evolutionary algorithms to solve each subcomponent.
However, the key insight was that the subcomponents could only be evaluated with each other, and therefore their fitness is determined by how well they solve the complete problem when combined with some representative solutions to the other subproblems.
There has been decades of work \citep{ma2018survey, pan2020effective, gong2015distributed, goh2008competitive, antonio2017coevolutionary,hancer2025survey, wiegand2004analysis, garcia2005cooperative, wiegand2001empirical, byrski2015evolutionary} building on this idea and endeavoring to solve the many issues that arise with implementing it: how to decompose the problem, how to decide on individual subcomponents to combine, how to evaluate each subcomponent, as well as the many other typical concerns of evolutionary algorithms such as mutation regime, etc.

However, this classic framing is a clear case of non-symbiotic co-evolution, because each individual subcomponent is only guaranteed to be evaluated with another individual subcomponent once.
This structure is analogous to an individual hummingbird visiting an individual flower: they both benefit if the hummingbird's beak happens to perfectly fit the flower's shape, but the hummingbird's fitness is based on how well it matches the flowers on average (taking into account the hummingbird's ability to choose flowers of course) and the flower's fitness is based on how well it matches the visiting pollinators on average.
A life-long relationship between individuals of different species, i.e., a symbiotic relationship, guarantees that two individuals will continuously interact with each other, regardless of whether the relationship is parasitic or mutualistic. These relationships can be highly asymmetric, with one organism being considered a \textit{host} and the other the \textit{symbiont}, in which case the ``life-long'' relationship is only in terms of the shorter of the two lifespans. These relationships are also frequently many-to-one (such as the many microbiomes that maintain the health of nearly all megafauna), or even many-to-many (such as the complex mycorrhizal fungi networks, where fungi can connect to several plant partners at once). 
Regardless of the specifics of the relationship, individuals in these symbiotic relationships are not being selected on their ability to match the average, but instead their ability to match their particular partner(s). These aspects of selection based on matching a specific partner and strong asymmetry in lifespans differ from how most co-evolutionary algorithms are implemented and could lead to a larger diversity of successful phenotypes.

Mutualistic symbiosis frequently also involves vertical transmission, where the host offspring is infected with the symbiont's offspring near birth. Vertical transmission is most often in the context of a host/symbiont system, i.e. when there is a strong asymmetry in the size and lifespan of the partners. This mode of transmission enables the offspring of the partners to continue the same mutualistic relationship.
Thus, the fitness of host and symbiont are linked and therefore the symbiont has strong selective pressure to continue to benefit the host.
The alternative to vertical transmission is horizontal transmission, when the symbiont reproduces by infecting other hosts in the population that are not necessarily related to the current host.
This route of transmission decouples the host's and endosymbiont's respective fitness and therefore can select for parasitism.
Therefore, the mode of transmission enabled by a co-evolutionary algorithm should take into account the form of relationship that is desired. If the co-evolutionary algorithm contains co-evolving tests and solutions, the tests can be viewed as parasites of the solutions and horizontal transmission should be used in the algorithm. 
Conversely, if the co-evolutionary algorithm consists of co-evolving subcomponents that should work mutualistically, vertical transmission may lead to better results.

Further, mutualism and parasitism are on a spectrum of antagonistic versus cooperative interactions and it is possible and likely that co-evolving relationships (both symbiotic and non-symbiotic) will evolve along that spectrum, leading to unexpected and perhaps undesirable effects if not controlled.
For example, if the goal is a population of co-evolving subcomponents that must partner together for a complete solution, propagating the subcomponents separately from their partners can lead to a break down of the mutualistic relationship and select for subcomponents that succeed at their partners' expense. This would likely then lead to suboptimal results overall and could at times explain why such an EC system fails to solve a problem.
	
In addition to the mutualism-parasitism spectrum, there are many ways in which symbiotic partners can impact each other, which ultimately all impact their fitness, but may have differing selective effects.
The typical co-evolutionary algorithm setup is best labeled as a cross-feeding dynamic, where no subcomponent is able to produce all necessary nutrients and so they must rely on each other for survival.
Specifically, each subcomponent can only solve part of the problem (i.e. a single necessary nutrient), and therefore must rely on the other subcomponents to produce a complete solution (i.e. the other necessary nutrients).
However, symbiotic partners can provide other types of benefits, such as protection against environmental stressors or pathogens \citep{Trivedi2020} and overall improved health \citep{Stengel2022}.
These other kinds of interactions may lead to differing selective effects due to the differing amounts of consistency (presence or absence of stressor) and directness (improved health) \citep{buckling2002role, benard2020stress, apprill2020role}, which may be beneficial to a particular problem domain.

Further, many other kinds of complications can occur.
For example, it has previously been shown that parasites can increase host behavioral complexity through the Population-Genetic Memory Hypothesis and Red Queen dynamics \citep{Ladle1992,Zaman2014}.
Mutualists, however, are hypothesized to have mixed effects.
They can increase host behavioral complexity through the major evolutionary transition of symbiogenesis, where the symbiont becomes so dependent on the host that it is no longer considered a separate organism \citep{CavalierSmith2013}; this biological phenomenon is somewhat analogous to the dynamics of Tangled Program Graphs \citep{Smith2021}.
Mutualists can also decrease host complexity by constraining host evolution, causing co-extinction events.
Mutualist constraint could also enable division of labor, decreasing the need for the host to perform complex tasks. While the host would then be less complex, such scenarios can actually improve the system's overall problem-solving capability \citep{Uchiumi2020}.
Endosymbiosis (where the symbiont lives inside of the host) has also been found to lead to genome degradation of the symbiont as it becomes reliant on the host, potentially decreasing the symbiont's behavioral complexity \citep{Andersson1999}.
It has also been hypothesized that the presence of both mutualistic and parasitic interactions in the same system could increase evolved host complexity beyond what the host would reach with only one type of interaction due to the more complex selection pressures from multiple interactions \citep{Willkens2022}.
Finally, when symbionts evolve along the parasitism-mutualism spectrum, the presence of spatial structure can influence how parasitic or mutualistic they become \citep{Vostinar2019}, in particular making parasitism more likely at intermediate vertical transmission rates. Further, their evolutionary history, such as the degree of specialization of the symbiont, can influence the future state of the system, as in the Co-Opted Antagonist Hypothesis, where a parasite is co-opted to become a mutualist due to increasing reliance on the host for survival\citep{Johnson2021}.

In the biological domain, host-symbiont interactions are strongly influenced by tangled mechanistic contingencies and often highly entrenched; we have so far not discovered a model system capable of every type of symbiotic interaction with which to perform controlled experiments (and it's unlikely such a system exists).
This limitation means that biological theory lacks a unified understanding of how each of these symbiotic interactions does impact the co-evolutionary dynamics of the system.
However, co-evolutionary algorithms do not have such a limitation, and so further development of these kinds of interactions within the EC field could benefit the development of such a framework.

Overall, symbiosis tends to complicate selective dynamics due to the highly individualized fitness landscapes that arise from a continuous interaction with a single and specific partner and so is not suited to all types of problems.
However, it has the potential to align and stabilize the selective pressure on subcomponents in a co-evolutionary algorithm, leading to more fruitful and diverse whole solutions.

% \putbib[biological-perspectives-coevolution]
    
% \end{bibunit}
